\subsection{Multi-Target Tracking Module}
\label{subsec:tracking}

At each sensing update $t$, the sensing module produces a set of candidates for the detected objects, each with a 3D representative position and an accumulated power.
However, when there is occlusion or clutter, the sensing module may produce ghost detections or miss detections.
To address this issue and achieve temporal consistency, the multi-target tracking module converts these detections into temporally consistent tracked targets.
Then, these traced targets are passed to the beamforming module.

Let the detection set at sensing update $t$ be
\begin{equation}
    \mathcal{D}_t =
    \{(\mathbf{z}_{t,j}, P_{t,j})\}_{j=1}^{N_t},
\end{equation}
where $\mathbf{z}_{t,j}\in\mathbb{R}^{3}$ is the representative position of detection $j$ and $P_{t,j}$ is its total power.
The sensing module uses power to rank candidate detections, while the tracking module uses position to match detections over time.

Each active tracked target is represented by a constant-velocity Kalman filter with state
\begin{equation}
    \mathbf{x}_{t,i}
    =
    [x_{t,i}, y_{t,i}, z_{t,i},
    v_{x,t,i}, v_{y,t,i}, v_{z,t,i}]^{\!\top}.
\end{equation},
where the first three components represent the positions and the last three components represent the velocities of the $i$-th tracked target.

At each sensing update$t$, the next possible state of each active tracked target is predicted as
\begin{align}
    \mathbf{x}_{t|t-1,i} & = F\mathbf{x}_{t-1|t-1,i},        \\
    P_{t|t-1,i}          & = F P_{t-1|t-1,i} F^{\!\top} + Q,
\end{align}
where $Q$ is the process-noise covariance and
\begin{equation}
    F =
    \begin{bmatrix}
        I_3                    & \Delta t I_3 \\
        \mathbf{0}_{3\times 3} & I_3
    \end{bmatrix}.
\end{equation}
$\Delta t$ corresponds to one sensing update.

After prediction, the tracking algorithm matches the predicted tracked targets to the detections in the current sensing update.
This matching process is known as data association in multi-target tracking.
For each predicted tracked target $i$ and detection $j$, the matching cost is the Euclidean distance
\begin{equation}
    c_{ij}
    =
    \left\|
    H\mathbf{x}_{t|t-1,i} - \mathbf{z}_{t,j}
    \right\|_2,
    \qquad
    H = [I_3 \;\; \mathbf{0}_{3\times 3}].
\end{equation}

Pairs with $c_{ij}>c_{\max}$ are considered unlikely and are assigned a large penalty.
The Hungarian algorithm~\cite{ref:hungarian} then selects the set of matches with the minimum total cost:
\begin{equation}
    \pi^\star
    =
    \arg\min_{\pi}
    \sum_{(i,j)\in\pi} c_{ij}.
\end{equation}
This joint matching avoids greedy decisions when multiple targets and detections are close to each other.

Each accepted tracked target--detection pair is corrected using the Kalman update.
The measurement residual and Kalman gain are
\begin{align}
    \mathbf{y}_{t,i} & =
    \mathbf{z}_{t,j} - H\mathbf{x}_{t|t-1,i}, \\
    S_{t,i}          & =
    H P_{t|t-1,i} H^{\!\top} + R,             \\
    K_{t,i}          & =
    P_{t|t-1,i} H^{\!\top} S_{t,i}^{-1},
\end{align}
where $R$ is the measurement-noise covariance.

The corrected state and covariance are then
\begin{align}
    \mathbf{x}_{t|t,i} & =
    \mathbf{x}_{t|t-1,i} + K_{t,i}\mathbf{y}_{t,i}, \\
    P_{t|t,i}          & =
    (I - K_{t,i}H)P_{t|t-1,i}.
\end{align}
After a successful update, the age of the tracked target is increased and its missed-frame counter is reset.

Unmatched predicted targets are kept temporarily by increasing their missed-frame counter, which allows the tracker to bridge short sensing gaps.
They are removed once this counter reaches the maximum allowed value $M_{\max}$.
For detections that are not matched to any of the existing tracked targets, a new tracked target is initialized with the detection's position and zero velocity:
\begin{equation}
    \mathbf{x}_{0}
    =
    [z_x, z_y, z_z, 0, 0, 0]^{\!\top}.
\end{equation}
The initial position is also stored as $\mathbf{p}_{\mathrm{start}}$ for later confirmation.

The tracker exports only confirmed moving tracked targets.
A tracked target is confirmed when it has survived for at least $A_{\min}$ frames and its displacement from the starting position exceeds $d_{\min}$:
\begin{equation}
    \left\|
    H\mathbf{x}_{t,i} - \mathbf{p}_{\mathrm{start},i}
    \right\|_2
    >
    d_{\min}.
\end{equation}
This rule reduces short-lived detections and static residual clutter.
In this work, we set  $M_{\max}=3$, $A_{\min}=3$, and $d_{\min}$ is configurable at runtime.
The output of this tracking module is the list of confirmed 3D target positions used by the sensing-informed beamforming stage.